% Pro Sex. Roscio Amerino --- headnote
% pro-roscio-amerino, 80 BC

\textit{Delivered at Rome in 80 BC, when Cicero was twenty-six and a
\textsl{novus homo} only three years past his first surviving speech.
The case is a defence of Sextus Roscius the younger, of Ameria, against
a charge of parricide; the real target of Cicero's attack is Lucius
Cornelius Chrysogonus, the most powerful of Sulla's freedmen, who had
acquired the murdered father's estates at a corrupt sale and now wished
the son condemned in order to silence the only witness. The speech
secured Roscius's acquittal and made Cicero's reputation: it is the
first work in which the rhetorical and political voice that defines the
later corpus is fully present, and the first in which Cicero attacks a
named representative of the Sullan order. He does so with elaborate
deference to Sulla himself, naming Chrysogonus as the architect of every
crime while protesting that the dictator's victory was just and willed
by the gods.}
